\documentclass[12pt]{article}
\usepackage[utf8]{inputenc}
\usepackage[T1]{fontenc}
\usepackage[margin=1in]{geometry}
\usepackage{setspace}
\usepackage{xcolor}
\usepackage[colorlinks=true, urlcolor=blue!60!black, linkcolor=blue!60!black]{hyperref}
\setlength{\parindent}{2em}
\setlength{\parskip}{1.25em}
\renewcommand{\baselinestretch}{1.0}
\begin{document}
\begin{center}
{\LARGE \textit{California v. Texas}, 593 U.S. 659 (2021)\par}
\vspace{0.8em}
{\large Case Brief --- Four-Page Edition\par}
\end{center}
\vspace{1em}

\section{Memory Jogger}

The Court dismissed the third major constitutional challenge to the Affordable Care Act, 7--2, holding that neither individual nor state plaintiffs had Article III standing to challenge an individual mandate whose penalty Congress had reduced to \$0 (\href{https://www.courtlistener.com/opinion/4892571/california-v-texas/}{593 U.S.\ 659}).

\section{Detailed Case Facts}

The ACA's minimum essential coverage provision, the ``individual mandate,'' 26 U.S.C.\ \S\ 5000A(a), told individuals to maintain coverage, and an income-scaled ``shared responsibility payment'' enforced it. In \textit{NFIB v.\ Sebelius} (2012), five Justices rejected Commerce Clause authority for the mandate, but the Court upheld \S\ 5000A under the Taxing Power because the exaction raised revenue. In 2017, after repeal efforts failed, the Tax Cuts and Jobs Act set the payment at \$0 beginning in 2019 while leaving the coverage language in place, and the IRS dropped the coverage check-box from Form 1040.

In 2018, Texas and 17 other States, later joined by two Texas residents, Neill Hurley and John Nantz, sued federal officials. They argued that a \$0 penalty raises no revenue, so \S\ 5000A(a) could no longer rest on the Taxing Power, and that it was inseverable from the rest of the ACA. Because the Trump Administration declined to defend the mandate, California, 15 other States, and the District of Columbia intervened to defend it, later joined by more States and the House of Representatives.

\section{Procedural History}

The Northern District of Texas held that the individual plaintiffs had standing, that \S\ 5000A(a) was unconstitutional, and that it was inseverable from the entire ACA, staying its judgment pending appeal. A divided Fifth Circuit found standing for both sets of plaintiffs and affirmed the constitutional ruling but remanded for a provision-by-provision severability analysis. The Supreme Court granted cross-petitions, heard argument on November 10, 2020, and on June 17, 2021, reversed on standing, vacated, and remanded with instructions to dismiss.

\section{Judicial Votes}

Dismissed for lack of standing, \textbf{7--2}.
\begin{itemize}
\item \textbf{Majority:} Justice Breyer, joined by Chief Justice Roberts and Justices Thomas, Sotomayor, Kagan, Kavanaugh, and Barrett.
\item \textbf{Concurring:} Justice Thomas, joining the Court's opinion in full.
\item \textbf{Dissenting:} Justice Alito, joined by Justice Gorsuch.
\end{itemize}

\section{Holding}

Neither the individual nor the state plaintiffs showed a past or future injury fairly traceable to any federal official's actual or threatened enforcement of \S\ 5000A(a), so the Court did not reach its constitutionality or severability. Standing to challenge a statute requires an injury traceable to enforcement of the specific provision attacked and redressable by relief against specific parties; unenforceable statutory text carrying no penalty cannot by itself supply traceability or redressability. Because the Court rejected ``standing-through-inseverability'' on forfeiture grounds, it left open whether that theory could succeed if timely preserved.

\section{Analysis of Opinions}

\textbf{Majority (Breyer).} Applying \textit{Lujan}'s injury, traceability, and redressability elements, Breyer held that even if premiums were a pocketbook injury, the individuals' costs were not traceable to enforcement: with a \$0 penalty, no federal officer could enforce \S\ 5000A(a) against them. Remedies operate against officials, not abstract rules, and the Declaratory Judgment Act confers no jurisdiction of its own, so a declaration would be an advisory opinion. The States' indirect theory, that the mandate drove residents into Medicaid, CHIP, and state employee plans, failed under \textit{Department of Commerce v.\ New York} and \textit{Clapper}: it rested on third-party choices without proof, only four of 21 declarations tied costs to the mandate (all from the penalty era), and a single CBO line did not show enrollment in state plans. Their direct costs, such as reporting under \S\S\ 6055 and 6056 and the employer mandate, flowed from separate provisions they did not challenge.

\textbf{Thomas (concurring).} Thomas acknowledged the force of Alito's critique of \textit{NFIB} and \textit{King v.\ Burwell} and the defenders' reversal on whether the mandate was essential. He suggested standing-through-inseverability might work in a future case if framed as a merits question of statutory construction, but declined to apply it because it was not raised below, not clearly briefed, and supported only by ``drive-by jurisdictional rulings.''

\textbf{Alito (dissenting).} Alito called the decision the third chapter of an ACA trilogy in which the Court rescued the statute through strained reasoning. The States suffered concrete costs, including reporting expenses and coverage of adult children up to age 26, backed by federal fines. Misreading \textit{Allen v.\ Wright}, the majority required tracing to the unconstitutional provision rather than to the challenged conduct, contrary to cases from \textit{Seila Law} back to the \textit{Trade-Mark Cases}. On the merits, a \$0 penalty raises no revenue, ``shall'' is mandatory, and Congress's finding that the mandate was ``essential'' made it inseverable from the provisions burdening the States.

\section{Examples: Future Applications}

\begin{enumerate}
\item \textbf{Same-Side (No Standing): Zero-Penalty Composting Directive.} A homeowner who buys bins to obey a \$0-penalty federal ``shall compost'' statute has no injury traceable to enforcement; a declaration would be advisory.
\item \textbf{Same-Side (No Standing): Speculative State Rebate Costs.} States claiming higher rebate outlays because residents obey an unenforceable smoke-detector directive offer only a speculative chain under \textit{Clapper}.
\item \textbf{Opposite-Side (Standing): Enforceable Assessment.} A \$200 IRS assessment collected by refund offset gives an official real enforcement authority, so an injunction against collection would redress the injury.
\item \textbf{Opposite-Side (Standing): Collateral Debarment.} A \$0 fine paired with automatic exclusion from federal contracting supplies a concrete enforcement mechanism a court can enjoin.
\item \textbf{Fence-Sitter: Preserved Inseverability Theory.} A utility burdened by an enforceable reporting rule, tied by an express inseverability clause to an unenforceable \$0 directive, pleads standing-through-inseverability from the start. Both the majority and Thomas left that question open.
\end{enumerate}

\section{Critique}

\textbf{Doctrinal ambiguity.} The \textit{Harvard Law Review} observed that no opinion engaged \textit{Massachusetts v.\ EPA}'s ``special solicitude'' for state plaintiffs, leaving lower courts to guess whether it had been quietly abandoned. Critics also debated the line between standing and merits: Alito's objection exposes tension with earlier inseverability cases, while defenders warn that letting any regulated party attack unrelated, unenforceable clauses would invite oversight of symbolic legislation.

\textbf{Ideological views.} Progressives welcomed a result that safeguarded coverage for millions but noted that it joined \textit{TransUnion} and \textit{Carney v.\ Adams} in tightening standing in ways that may later hinder challenges to deregulation. Textualists split: some praised the Court for refusing advisory declarations, while others shared Alito's view that standing rules were applied selectively to insulate the ACA.

\section{Key Quotations}

\begin{quote}
``To find standing here to attack an unenforceable statutory provision would allow a federal court to issue what would amount to `an advisory opinion without the possibility of any judicial relief.' ''
\end{quote}
\noindent --- Justice Breyer, for the Court.

\begin{quote}
``The provision is unenforceable. There is no one, and nothing, to enjoin.''
\end{quote}
\noindent --- Justice Breyer, for the Court.

\begin{quote}
``A penalty is a tax. The United States is a State. And 18 States who bear costly burdens under the ACA cannot even get a foot in the door to raise a constitutional challenge.''
\end{quote}
\noindent --- Justice Alito, dissenting.

\end{document}
